\documentclass[11pt,a4paper]{article}

\usepackage[utf8]{inputenc}
\usepackage[T1]{fontenc}
\usepackage{lmodern}
\usepackage{amsmath,amssymb}

% Prevent graphicx from loading epstopdf-base (often missing in minimal
% TeX Live installs).  Declaring a fake version satisfies the dependency.
\makeatletter
\@namedef{ver@epstopdf-base.sty}{2024/01/01 stub}
\makeatother

\usepackage{graphicx}
\usepackage{geometry}
\usepackage{hyperref}
\usepackage{booktabs}
\usepackage{xcolor}

\geometry{margin=2.5cm}

\setlength{\parindent}{0pt}
\setlength{\parskip}{6pt plus 2pt minus 1pt}

\hypersetup{
    colorlinks=true,
    linkcolor=blue!60!black,
    urlcolor=blue!60!black,
    citecolor=green!50!black,
}

\pagestyle{headings}

\title{\textbf{RICH Regression Test Suite} \\ \large Documentation and Results}
\author{Auto-generated by \texttt{generate\_test\_report.py}}
\date{\today}

\begin{document}

\maketitle
\tableofcontents
\newpage


\section{Verification and Validation}

\subsection{Background and Definitions}

Verification and Validation (V\&V) is the primary means by which confidence in
computational simulations is established \cite{roache1998,oberkampf2010}.
Although the two terms are frequently paired, they address fundamentally
distinct questions:

\begin{itemize}
  \item \textbf{Verification} asks: \emph{``Are we solving the equations
    correctly?''}  It is a \emph{mathematics} activity that checks whether the
    numerical implementation faithfully reproduces the intended mathematical
    model.  Verification is typically decomposed into \emph{code verification}
    (detecting programming errors by comparing against known analytical or
    manufactured solutions) and \emph{solution verification} (estimating
    numerical discretisation errors through convergence studies)
    \cite{roache1998,salari2000}.
  \item \textbf{Validation} asks: \emph{``Are we solving the correct
    equations?''}  It is a \emph{physics} activity that compares simulation
    predictions with experimental or observational data to assess whether the
    mathematical model adequately represents the physical phenomena of interest
    \cite{oberkampf2010}.
\end{itemize}

The distinction was formalised by the American Institute of Aeronautics and
Astronautics (AIAA) \cite{aiaa1998} and later expanded into comprehensive
guidelines by the American Society of Mechanical Engineers (ASME)
\cite{asme2006}.  Together, V\&V provides a systematic framework for
quantifying and reducing uncertainty in simulation results, an essential step
before codes are used in predictive settings such as high-energy-density physics
(HEDP), inertial confinement fusion (ICF), astrophysics, and engineering
design.

\subsection{The V\&V Process}

A robust V\&V programme typically proceeds through several stages
\cite{oberkampf2010}:

\begin{enumerate}
  \item \textbf{Code verification with exact solutions.}
    Problems that possess known closed-form or semi-analytical solutions---such
    as the Sod shock tube \cite{sod1978}, the Sedov--Taylor blast wave
    \cite{sedov1959,taylor1950}, or Lane--Emden polytropic equilibria
    \cite{chandrasekhar1939}---are used to confirm that the code reproduces the
    correct answer.  Error norms are monitored and, where feasible, the order of
    convergence is verified against the theoretical rate of the discretisation.

  \item \textbf{Code verification with manufactured solutions.}
    When analytical solutions are unavailable, the Method of Manufactured
    Solutions (MMS) \cite{roache2002,salari2000} introduces a prescribed
    solution and derives the source terms required to sustain it.  The code is
    then run with those source terms, and the discretisation error is checked
    for the expected convergence rate.

  \item \textbf{Solution verification.}
    Even after code correctness is established, discretisation and iteration
    errors remain.  Grid-convergence studies (e.g.\ Richardson extrapolation
    \cite{richardson1911,roache1998}) and time-step refinement are used to
    estimate and bound these errors in production calculations.

  \item \textbf{Validation against experimental or observational data.}
    Finally, simulations are compared with real-world measurements.  Agreement
    gives confidence; discrepancies drive model improvement.  Validation is
    inherently problem-specific and requires careful attention to experimental
    uncertainties \cite{oberkampf2010}.
\end{enumerate}

In practice, \emph{regression testing}---re-running a fixed set of verification
problems after every code change---is the mechanism that ensures
previously passing tests continue to pass.  Regression testing does not replace
formal convergence studies or validation campaigns, but it provides a
continuous, automated safety net that catches inadvertent errors early in the
development cycle.

\subsection{Role of the RICH Regression Test Suite}

The test suite documented in this report serves the \emph{code verification}
and \emph{regression testing} aspects of V\&V for the RICH code.  Each
regression test exercises a specific subset of the code's physics
capabilities---compressible hydrodynamics, flux-limited radiation diffusion,
multigroup radiation transport, Compton scattering, Newtonian self-gravity, and
adaptive mesh refinement (AMR)---and compares the results against known
analytical solutions, published reference data, or strict conservation
invariants.

The suite is designed with the following principles:

\begin{enumerate}
  \item \textbf{Correctness.}  Detect coding errors, regressions, or
    platform-dependent numerical issues as early as possible by comparing
    simulation output to known solutions.

  \item \textbf{Reproducibility.}  Every test is fully automated, deterministic,
    and produces machine-readable pass/fail signals (exit codes) so that the
    suite can be run routinely in continuous-integration or pre-merge workflows.

  \item \textbf{Coverage.}  The suite spans serial and MPI execution, Eulerian
    and Lagrangian mesh motion, single- and multi-group radiation, and multiple
    spatial dimensions, providing broad coverage of the code's feature space.

  \item \textbf{Documentation.}  For each test, this report records the problem
    definition (initial and boundary conditions, mesh movement, material
    properties), the quantitative pass criteria with explicit thresholds, and
    the achieved numerical results, forming a living document that evolves with
    the code.
\end{enumerate}

The individual tests are summarised in the table below and described in detail
in the sections that follow.  For each test we provide references to the
underlying analytical solutions or prior work on which the problem setup is
based.

\newpage

\section{Summary}
\label{sec:summary}

\begin{table}[htbp]
\centering
\caption{Overview of all regression tests.}
\label{tab:summary}
\begin{tabular}{l l l l l}
\toprule
\textbf{Test} & \textbf{Mode} & \textbf{CPUs} & \textbf{Mesh} & \textbf{Plots} \\
\midrule
Sod 1D & Serial & 1 & Eulerian & Yes \\
Sedov 3D & MPI & 128 & Lagrangian & Yes \\
Till Compton & Serial & 1 & Lagrangian & Yes \\
AMR Random & Serial + MPI & 1 / 64 & Lagrangian + AMR & No \\
Voronoi Volume & Serial + MPI & 1 / 64 & Static & No \\
Lane--Emden & MPI & 64 & Lagrangian & Yes \\
Mach 2 Gray & MPI & 8 & Eulerian & Yes \\
Mach 2 Multigroup & MPI & 8 & Eulerian & Yes \\
Marshak Wave 1 & Serial & 1 & Eulerian & Yes \\
Marshak Wave 2 & Serial & 1 & Eulerian & Yes \\
Marshak Wave 3 & Serial & 1 & Eulerian & Yes \\
Marshak Wave 4 & Serial & 1 & Eulerian & Yes \\
Gresho Vortex (Euler) & Serial & 1 & Eulerian & Yes \\
Gresho Vortex (Lagrangian) & MPI & 8 & Lagrangian & Yes \\
Spherical LSQ Gradient & Serial & 1 & Static & No \\
Spherical Collapse & MPI & 64 & Eulerian & Yes \\
Rayleigh--Taylor & MPI & 128 & Lagrangian & Yes \\
Gray Free--Free Suite & MPI & 4--16 & Eulerian & Yes \\
Multigroup Free--Free Suite & MPI & 4--16 & Eulerian & Yes \\
\bottomrule
\end{tabular}
\end{table}

\newpage

\section{Sod Shock Tube (1D)}
\label{sec:sod_1d}

\subsection*{Description}
The Sod shock tube \cite{sod1978} is a classical one-dimensional Riemann problem that produces a left-moving rarefaction wave, a contact discontinuity, and a right-moving shock. It is one of the most widely used verification benchmarks for compressible hydrodynamics codes \cite{toro2009}.

\textbf{Code and physics aspects verified:}
\begin{itemize}
  \item \textbf{Riemann solver (HLLC):} The test directly exercises the approximate Riemann solver that computes inter-cell fluxes. Correct shock speed, contact velocity, and rarefaction fan structure all depend on the solver functioning properly.
  \item \textbf{Spatial reconstruction (PLM):} The piecewise-linear reconstruction with slope limiting determines the order of accuracy and controls spurious oscillations near discontinuities. A faulty limiter would produce either excessive diffusion or non-physical overshoots at the shock and contact.
  \item \textbf{Equation of state:} The ideal-gas EOS ($\gamma$-law) is tested through the pressure--density--energy coupling across all three wave families.
  \item \textbf{Time integration and CFL condition:} The Courant time-step limiter and the conservative update must correctly advance the solution without introducing instability.
\end{itemize}

\textbf{Importance:} The Sod problem is the most fundamental hydrodynamics test. Because it has an exact analytical solution, it provides an unambiguous correctness check. A code that fails this test cannot be trusted for any compressible-flow simulation. Its one-dimensional, Eulerian nature also makes it a fast, inexpensive smoke test that can be run on every commit.

\subsection*{Initial Conditions}
The domain is $x \in [0,\,1]$ with 400 uniformly spaced cells. The initial discontinuity is at $x = 0.5$:
\begin{itemize}
  \item \textbf{Left state} ($x < 0.5$): $\rho = 1$, $P = 1$, $u = 0$.
  \item \textbf{Right state} ($x \ge 0.5$): $\rho = 0.125$, $P = 0.1$, $u = 0$.
\end{itemize}
The adiabatic index is $\gamma = 1.4$. The simulation runs to $t = 0.2$.

\subsection*{Boundary Conditions}
Rigid (reflective) walls at both ends of the domain.

\subsection*{Mesh Movement}
Eulerian (fixed mesh).

\subsection*{Execution}
Serial, 1~CPU, direct execution.

\subsection*{Pass Criteria}
The numerical density and pressure profiles are compared to the exact Riemann solution (computed via the \texttt{enrs.py} solver). The goodness-of-fit (GOF) metric must satisfy:
\begin{itemize}
  \item Density GOF $\le 0.02$.
  \item Pressure GOF $\le 0.02$.
\end{itemize}

\subsection*{Achieved Results}
\begin{table}[htbp]
\centering
\begin{tabular}{l r r c}
\toprule
\textbf{Metric} & \textbf{Achieved} & \textbf{Threshold} & \textbf{Status} \\
\midrule
Density GOF & 1.14136222e-02 & 2.00000000e-02 & {\color{green!60!black}PASS} \\
Pressure GOF & 1.07895027e-02 & 2.00000000e-02 & {\color{green!60!black}PASS} \\
\bottomrule
\end{tabular}
\end{table}

\subsection*{Plots}
\textit{Plots not available --- run the test suite and the plot generator first.}

\newpage

\section{Sedov--Taylor Blast Wave (3D, MPI)}
\label{sec:sedov_3d_mpi}

\subsection*{Description}
The Sedov--Taylor blast wave \cite{sedov1959,taylor1950} is a point-explosion problem with a self-similar analytical solution obtained by integrating an ODE \cite{kamm2007}. It is a standard verification problem for multi-dimensional Lagrangian and ALE hydrodynamics.

\textbf{Code and physics aspects verified:}
\begin{itemize}
  \item \textbf{3D Voronoi tessellation:} The unstructured Voronoi mesh is generated from $\sim\!5\times10^6$ random seed points and must correctly compute cell volumes, face areas, and geometric connectivity in three dimensions.
  \item \textbf{Lagrangian mesh motion:} Cell generators move with the local fluid velocity, with periodic cell-rounding corrections to maintain mesh quality. This tests the full mesh-motion pipeline including velocity interpolation, generator advection, and incremental re-tessellation.
  \item \textbf{Spherical symmetry on an unstructured mesh:} A strong spherical shock must propagate symmetrically despite the inherently irregular Voronoi geometry. This stresses the isotropy of the numerical scheme.
  \item \textbf{MPI domain decomposition:} The problem is distributed across 128 MPI ranks. Correct ghost-cell exchange, parallel tessellation, and flux communication are all exercised.
  \item \textbf{Strong-shock hydrodynamics:} The Sedov blast wave involves an infinite-strength shock with a compression ratio of $(\gamma+1)/(\gamma-1) = 4$ (for $\gamma = 5/3$). Correctly capturing the peak density and the post-shock profile is demanding.
\end{itemize}

\textbf{Importance:} This is the primary multi-dimensional, Lagrangian, parallel verification test. It simultaneously checks the geometric machinery (Voronoi), the physics (strong-shock hydrodynamics), and the parallel infrastructure (MPI). Failures here can indicate bugs ranging from tessellation errors to conservation violations across processor boundaries.

\subsection*{Initial Conditions}
The domain is a cube $[-1,\,1]^3$ filled with a Voronoi tessellation from $\sim\!5\times10^6$ random points. Cells with centroid $r < 0.2$ receive high energy:
\begin{itemize}
  \item \textbf{Inner region} ($r < 0.2$): $\rho = 1$, $e_{\mathrm{int}} = 10^5$.
  \item \textbf{Outer region} ($r \ge 0.2$): $\rho = 1$, $e_{\mathrm{int}} = 0.1$.
\end{itemize}
The adiabatic index is $\gamma = 5/3$.

\subsection*{Boundary Conditions}
Rigid (reflective) walls on all six faces of the cube.

\subsection*{Mesh Movement}
Lagrangian (cells move with the fluid velocity), with cell rounding.

\subsection*{Execution}
MPI, 128~CPUs, submitted via SLURM (partition \texttt{bigrun}, exclusive).

\subsection*{Pass Criteria}
Radially binned profiles of density, pressure, and velocity are compared to the Sedov--Taylor ODE solution (scaled from theory using $v_s = \tfrac{2}{5}\,R_s/t$, not normalized to numerical peaks). The volume-weighted relative $L_1$ error (weights $\propto r^2$) must satisfy:
\begin{itemize}
  \item Density: relative $L_1 \le 0.50$.
  \item Pressure: relative $L_1 \le 0.30$.
  \item Velocity: relative $L_1 \le 0.60$.
\end{itemize}

\subsection*{Achieved Results}
\begin{table}[htbp]
\centering
\begin{tabular}{l r r c}
\toprule
\textbf{Metric} & \textbf{Achieved} & \textbf{Threshold} & \textbf{Status} \\
\midrule
Density rel. $L_1$ & 1.72107611e-01 & 5.00000000e-01 & {\color{green!60!black}PASS} \\
Pressure rel. $L_1$ & 4.97101881e-02 & 3.00000000e-01 & {\color{green!60!black}PASS} \\
Velocity rel. $L_1$ & 1.58704245e-02 & 6.00000000e-01 & {\color{green!60!black}PASS} \\
\bottomrule
\end{tabular}
\end{table}

\subsection*{Plots}
\textit{Plots not available --- run the test suite and the plot generator first.}

\newpage

\section{Till Compton Equilibration}
\label{sec:till_compton}

\subsection*{Description}
This test verifies the radiation--matter energy exchange via Compton scattering and absorption/emission in a uniform, static medium, following the problem setup of Till \cite{till2021}. The numerical results are compared against the IN-FBC (Inline Full-Boltzmann Compton) reference data from McGraw et al.\ \cite{mcgraw2023}. Starting from mismatched gas and radiation temperatures, the system should relax to thermal equilibrium.

\textbf{Code and physics aspects verified:}
\begin{itemize}
  \item \textbf{Compton scattering:} The dominant energy-exchange mechanism at keV temperatures. The test verifies that the Compton heating operator correctly transfers energy between the radiation field and the electrons, including the induced-scattering nonlinearity inherent to Bose--Einstein statistics.
  \item \textbf{Absorption and emission:} The Kramers free--free (bremsstrahlung) opacity couples each energy group to the material temperature. Correct group-by-group absorption/emission rates are essential for the system to relax to the right equilibrium.
  \item \textbf{Multigroup radiation transport:} The 32 logarithmically spaced energy groups must each evolve consistently, and the total radiation energy density must be conserved to within the implicit solver tolerance.
  \item \textbf{Implicit time integration:} The radiation--matter coupling is stiff (the Compton and absorption time scales are much shorter than the equilibration time). The implicit solver must remain stable and converge even for large time steps.
  \item \textbf{Energy conservation:} Total energy (radiation + material) must be conserved. Any leak indicates a bug in the coupling operator or the implicit solve.
\end{itemize}

\textbf{Importance:} Compton scattering is the primary radiation--matter coupling mechanism in high-energy-density physics (HEDP) and astrophysical applications. This test is the only one in the suite that exercises the Compton-scattering module in isolation (without hydrodynamic motion), making it essential for verifying the radiation physics independently of the hydro solver.

\subsection*{Initial Conditions}
A single Voronoi cell in a unit domain with density $\rho = 1\;\mathrm{g\,cm^{-3}}$ of hydrogen ($Z=1$). The initial gas temperature is $T_{\mathrm{gas}} = 1\;\mathrm{keV}$ and the initial radiation temperature is $T_{\mathrm{rad}} = 10\;\mathrm{keV}$. Multigroup radiation transport is used with 32 logarithmically spaced energy groups from $10^{-1}\;\mathrm{eV}$ to $10^{3}\;\mathrm{keV}$. Hydrodynamics is disabled.

The absorption opacity is the Kramers free--free (bremsstrahlung) formula:
\[
  \sigma_a(\nu,\,T) = 3.7\times10^{8}\;Z^{3}\,n_e\,n_i\;T^{-1/2}\,\bigl(1 - e^{-h\nu/k_BT}\bigr)\,\nu^{-3}\quad[\mathrm{cm^{-1}}],
\]
where $n_e = n_i = \rho\,N_A \approx 6.022\times10^{23}\;\mathrm{cm^{-3}}$ for fully ionized hydrogen at $\rho = 1\;\mathrm{g\,cm^{-3}}$. With $Z=1$ the leading constant evaluates to $3.7\times10^{8}\times(6.022\times10^{23})^{2} \approx 1.342\times10^{56}\;\mathrm{cm^{-1}\,K^{1/2}\,Hz^{3}}$. Scattering opacity is zero.

\subsection*{Boundary Conditions}
Rigid walls (irrelevant since hydrodynamics is off).

\subsection*{Mesh Movement}
Lagrangian, but the cell is effectively static (single cell, no flow).

\subsection*{Execution}
Serial, 1~CPU, direct execution. Built with \texttt{--energy\_groups\_num=32}.

\subsection*{Pass Criteria}
The simulation produces time histories of $T_{\mathrm{gas}}(t)$, $T_{\mathrm{rad}}(t)$, and the total specific energy $E_{\mathrm{tot}}(t) = e_{\mathrm{int}} + E_{\mathrm{rad}}/\rho$. The test passes when:
\begin{itemize}
  \item No NaN or Inf values appear in the output.
  \item The final gas and radiation temperatures agree within 1\%:
        $|T_{\mathrm{gas}} - T_{\mathrm{rad}}| / \max(T_{\mathrm{gas}},\,T_{\mathrm{rad}}) < 10^{-2}$.
  \item Total energy (thermal + radiation) is conserved:
        $|E_{\mathrm{tot}}^{\mathrm{final}} - E_{\mathrm{tot}}^{\mathrm{initial}}| / E_{\mathrm{tot}}^{\mathrm{initial}} < 10^{-8}$.
\end{itemize}

\subsection*{Achieved Results}
\begin{table}[htbp]
\centering
\begin{tabular}{l r r c}
\toprule
\textbf{Metric} & \textbf{Achieved} & \textbf{Threshold} & \textbf{Status} \\
\midrule
Final $T_{\mathrm{gas}}$ & 1.154070e+08 & --- & {\color{green!60!black}PASS} \\
Final $T_{\mathrm{rad}}$ & 1.154980e+08 & --- & {\color{green!60!black}PASS} \\
Relative difference & 7.878924e-04 & $ < 10^{-2}$ & {\color{green!60!black}PASS} \\
Energy conservation (rel.) & 0.000000e+00 & $ < 10^{-8}$ & {\color{green!60!black}PASS} \\
\bottomrule
\end{tabular}
\end{table}

\subsection*{Plots}
\textit{Plots not available --- run the test suite and the plot generator first.}

\newpage

\section{AMR Random Refine/Coarsen Consistency}
\label{sec:amr_random}

\subsection*{Description}
This test checks that adaptive mesh refinement (AMR) operations (random cell splitting and merging) preserve the hydrodynamic state of a uniform gas at rest. Any drift from the baseline indicates a conservation or interpolation bug in the AMR machinery.

\textbf{Code and physics aspects verified:}
\begin{itemize}
  \item \textbf{Cell refinement (splitting):} When a Voronoi cell is refined, the parent cell's conserved quantities (mass, momentum, energy) must be exactly partitioned among the children. The test verifies that the volume-weighted interpolation and the geometric insertion of new seed points are correct.
  \item \textbf{Cell coarsening (merging):} When cells are merged, the reverse operation must exactly recover the combined conserved quantities. Errors here would cause mass or energy to be created or destroyed.
  \item \textbf{Voronoi re-tessellation:} After adding or removing seed points, the tessellation must be rebuilt consistently. Cell volumes, face areas, and neighbour connectivity must all be updated without introducing geometric artefacts.
  \item \textbf{MPI consistency:} In the parallel variant, cells near processor boundaries may be refined or coarsened while their ghost copies reside on a different rank. The test checks that cross-rank AMR operations maintain consistency.
\end{itemize}

\textbf{Importance:} AMR is a critical capability for production simulations that need to resolve local features (shocks, contact discontinuities, collapse regions) without paying the cost of a globally fine mesh. Because AMR touches the deepest layers of the data structures---cell connectivity, conserved-variable storage, MPI communication patterns---even small bugs can cause silent conservation errors that accumulate over thousands of time steps. A uniform-gas test is the most stringent check: the answer must be exactly unchanged to machine precision.

\subsection*{Initial Conditions}
The domain is a cube $[-1,\,1]^3$ with a Voronoi tessellation from random points. The gas is uniform and at rest: $\rho = 1$, $e_{\mathrm{int}} = 2.5$, $\mathbf{v} = 0$.

\subsection*{Boundary Conditions}
Rigid (reflective) walls on all faces.

\subsection*{Mesh Movement}
Lagrangian with cell rounding. AMR randomly refines and coarsens cells during the simulation.

\subsection*{Execution}
Both serial and MPI. Serial: 1~CPU, direct execution. MPI: 64~CPUs via SLURM (partition \texttt{bigrun}, exclusive).

\subsection*{Pass Criteria}
After AMR operations, the maximum relative drift in density, internal energy, pressure, and velocity must stay small:
\begin{itemize}
  \item Serial: $\texttt{max\_drift} \le 10^{-8}$.
  \item MPI: $\texttt{max\_drift} \le 10^{-6}$.
\end{itemize}

\subsection*{Achieved Results}
\begin{table}[htbp]
\centering
\begin{tabular}{l r r c}
\toprule
\textbf{Metric} & \textbf{Achieved} & \textbf{Threshold} & \textbf{Status} \\
\midrule
Mode & mpi & --- & {\color{green!60!black}PASS} \\
Max drift & 1.1758199303102401e-08 & 9.9999999999999995e-07 & {\color{green!60!black}PASS} \\
\bottomrule
\end{tabular}
\end{table}

\subsection*{Plots}
\textit{This test does not produce comparison plots.}

\newpage

\section{Voronoi Volume Conservation}
\label{sec:voronoi_volume}

\subsection*{Description}
A pure-geometry test that verifies the Voronoi tessellation correctly partitions the domain. The sum of all cell volumes must match the known analytical box volume to machine precision.

\textbf{Code and physics aspects verified:}
\begin{itemize}
  \item \textbf{Voronoi geometry kernel:} The convex-hull and face-clipping algorithms that compute each cell's vertices, faces, and volume are tested end-to-end. A bug in vertex computation or face orientation would cause the total volume to deviate from the box volume.
  \item \textbf{Boundary handling:} Cells adjacent to the domain boundary are clipped against the bounding box. Correct clipping is essential for volume conservation and for imposing boundary conditions in production simulations.
  \item \textbf{Parallel tessellation:} In the MPI variant, each rank builds a local Voronoi diagram with ghost layers from neighbouring ranks. The sum of all local volumes must still equal the global box volume, verifying that no gaps or overlaps exist at processor boundaries.
\end{itemize}

\textbf{Importance:} The Voronoi tessellation is the geometric foundation of the entire code. Every flux calculation, gradient reconstruction, and cell-centred update depends on correct cell volumes and face areas. This test catches low-level geometric bugs (floating-point degeneracies, clipping errors, parallel inconsistencies) that would otherwise manifest as mysterious conservation violations in physics simulations. It is also extremely cheap to run, making it ideal as a first-line diagnostic.

\subsection*{Initial Conditions}
The domain is a unit cube $[0,\,1]^3$ filled with a Voronoi tessellation from random points. No hydrodynamics or time evolution is performed.

\subsection*{Boundary Conditions}
Rigid walls (domain boundary).

\subsection*{Mesh Movement}
Static (no time evolution).

\subsection*{Execution}
Both serial and MPI. Serial: 1~CPU, direct execution. MPI: 64~CPUs via SLURM (partition \texttt{bigrun}, exclusive).

\subsection*{Pass Criteria}
The relative volume error must be negligible:
\begin{itemize}
  \item $\texttt{rel\_error} = |V_{\mathrm{total}} - V_{\mathrm{box}}| / V_{\mathrm{box}} < 10^{-10}$.
\end{itemize}

\subsection*{Achieved Results}
\begin{table}[htbp]
\centering
\begin{tabular}{l r r c}
\toprule
\textbf{Metric} & \textbf{Achieved} & \textbf{Threshold} & \textbf{Status} \\
\midrule
Relative volume error & 0.0000000000000000e+00 & $ < 10^{-10}$ & {\color{green!60!black}PASS} \\
\bottomrule
\end{tabular}
\end{table}

\subsection*{Plots}
\textit{This test does not produce comparison plots.}

\newpage

\section{Lane--Emden Self-Gravity Equilibrium}
\label{sec:lane_self_gravity}

\subsection*{Description}
This test initializes a polytropic star in hydrostatic equilibrium according to the Lane--Emden equation \cite{chandrasekhar1939,lane1870} (polytropic index $n = 1.5$) and evolves it with self-gravity. The density profile should remain close to the analytical solution. This is a standard benchmark for codes coupling hydrodynamics and Newtonian gravity.

\textbf{Code and physics aspects verified:}
\begin{itemize}
  \item \textbf{Self-gravity solver:} The Newtonian gravitational potential and its gradient must be computed accurately enough that the gravitational acceleration balances the pressure gradient, maintaining hydrostatic equilibrium. Errors in the gravity solver would cause the star to spuriously expand or collapse.
  \item \textbf{Hydrostatic balance on an unstructured mesh:} Maintaining a static equilibrium with steep density gradients on a Voronoi mesh is challenging. The pressure-gradient reconstruction and the gravity-source coupling must be well-balanced to avoid generating spurious velocities.
  \item \textbf{Spherical geometry on a Cartesian domain:} The spherical star is embedded in a Cartesian bounding box. The test checks that the boundary between the star's interior and the low-density exterior is handled cleanly.
  \item \textbf{Conservation under long-time evolution:} The simulation runs for multiple dynamical times. Slow drifts in total energy or mass, or gradual distortion of the density profile, would indicate conservation or well-balancing defects.
\end{itemize}

\textbf{Importance:} Self-gravity is essential for astrophysical applications (star formation, stellar structure, compact-object mergers). The Lane--Emden equilibrium is one of the few configurations with an analytical density profile, making it invaluable for verifying the gravity module. It is also a sensitive test of the code's ability to maintain a static equilibrium---a prerequisite for reliable long-duration simulations of gravitationally bound systems.

\subsection*{Initial Conditions}
A sphere of mass $M = 2\times10^{33}\;\mathrm{g}$ and radius $R = 7\times10^{10}\;\mathrm{cm}$ is set up with the Lane--Emden density profile (read from tabulated data files \texttt{data/xsi32.txt} and \texttt{data/theta32.txt}). The gravitational constant is $G = 6.674\times10^{-8}\;\mathrm{cgs}$. The Voronoi mesh uses stratified random sampling inside the sphere.

\subsection*{Boundary Conditions}
Rigid walls at the outer boundary.

\subsection*{Mesh Movement}
Lagrangian with cell rounding. Gravity via a conservative force module.

\subsection*{Execution}
MPI, 64~CPUs, submitted via SLURM (partition \texttt{bigrun}, exclusive).

\subsection*{Pass Criteria}
The volume-weighted average density error vs.\ the analytical Lane--Emden profile must satisfy:
\begin{itemize}
  \item $|\texttt{final\_metric}| = \frac{\sum |\rho_{\mathrm{num}} - \rho_{\mathrm{analytic}}|\,V}{\sum V} < 4\times10^{-2}$.
\end{itemize}

\subsection*{Achieved Results}
\begin{table}[htbp]
\centering
\begin{tabular}{l r r c}
\toprule
\textbf{Metric} & \textbf{Achieved} & \textbf{Threshold} & \textbf{Status} \\
\midrule
|final\_metric| & 1.751990e-02 & 4.00e-02 & {\color{green!60!black}PASS} \\
\bottomrule
\end{tabular}
\end{table}

\subsection*{Plots}
\textit{Plots not available --- run the test suite and the plot generator first.}

\newpage

\section{Mach~2 Radiative Shock -- Gray Diffusion}
\label{sec:mach2_diffusion}

\subsection*{Description}
A Mach~2 radiative shock with gray (single-group) flux-limited diffusion \cite{lowrie2008,lowrie1999}. Radiative shocks exhibit a rich structure including a precursor region where radiation preheats the upstream material and a relaxation zone behind the shock \cite{mihalas1984,zel2002}. The numerical profiles are compared to the NLTE analytical solution for density, gas temperature, and radiation temperature.

\textbf{Code and physics aspects verified:}
\begin{itemize}
  \item \textbf{Radiation--hydrodynamics coupling:} This is the first test in the suite that couples radiation transport and hydrodynamics in a non-trivial flow. The radiation precursor heats the upstream gas before the shock arrives, modifying the shock structure. Getting both the precursor length and the post-shock relaxation zone correct requires accurate coupling between the radiation and hydro solvers.
  \item \textbf{Flux-limited diffusion:} The gray radiation diffusion equation with a flux limiter must transition smoothly between the diffusion regime (optically thick) and the free-streaming regime (optically thin). The flux limiter is critical for preventing superluminal radiation transport.
  \item \textbf{Implicit radiation solve:} The radiation diffusion equation is solved implicitly. The linear solver, matrix assembly, and convergence criteria are all exercised.
  \item \textbf{Non-equilibrium radiation:} The gas and radiation temperatures differ throughout the precursor and relaxation regions ($T_{\mathrm{gas}} \neq T_{\mathrm{rad}}$). The test verifies that the code correctly handles non-equilibrium conditions and evolves toward the correct downstream equilibrium.
  \item \textbf{Steady-state shock structure:} The simulation must reach a steady state in the shock frame, requiring correct inflow boundary conditions and long enough run time.
\end{itemize}

\textbf{Importance:} Radiative shocks are ubiquitous in HEDP experiments and astrophysical phenomena (accretion shocks, supernova remnants, stellar atmospheres). The Mach~2 Lowrie benchmark is one of the few radiation-hydrodynamics problems with a semi-analytical solution, making it a cornerstone verification test. Passing this test gives confidence that the code can simulate radiation-dominated flows where the radiation field significantly modifies the hydrodynamic structure.

\subsection*{Initial Conditions}
A 1D Cartesian mesh with 1024 cells spanning $x \in [-10^3,\;2\times10^3]\;\mathrm{cm}$. The left and right states are:
\begin{itemize}
  \item \textbf{Left (upstream):} $\rho = 5.459\times10^{-13}\;\mathrm{g/cm^3}$, $v = 2.355\times10^{5}\;\mathrm{cm/s}$, $T = 100\;\mathrm{K}$.
  \item \textbf{Right (downstream):} $\rho = 1.248\times10^{-12}\;\mathrm{g/cm^3}$, $v = 1.03\times10^{5}\;\mathrm{cm/s}$, $T = 207.8\;\mathrm{K}$.
\end{itemize}
Adiabatic index $\gamma = 5/3$; Rosseland opacity $\sigma_{\mathrm{Ross}} = 0.849\;\mathrm{cm^{-1}}$; absorption opacity $\sigma_{\mathrm{abs}} = 3.93\times10^{-5}\;\mathrm{cm^{-1}}$. The simulation time is $t = 0.01\;\mathrm{s}$.

\subsection*{Boundary Conditions}
Inflow boundaries at left and right.

\subsection*{Mesh Movement}
Eulerian (fixed mesh).

\subsection*{Execution}
MPI, 8~CPUs, submitted via SLURM (partition \texttt{bigrun}, exclusive).

\subsection*{Pass Criteria}
Profiles are compared to the NLTE analytical solution. The relative $L_1$ error must satisfy:
\begin{itemize}
  \item Density: relative $L_1 \le 0.025$.
  \item Temperature: relative $L_1 \le 0.025$.
\end{itemize}

\subsection*{Achieved Results}
\begin{table}[htbp]
\centering
\begin{tabular}{l r r c}
\toprule
\textbf{Metric} & \textbf{Achieved} & \textbf{Threshold} & \textbf{Status} \\
\midrule
Density rel. $L_1$ & 3.12941797e-03 & 2.50000000e-02 & {\color{green!60!black}PASS} \\
Temperature rel. $L_1$ & 6.17569973e-03 & 2.50000000e-02 & {\color{green!60!black}PASS} \\
$T_{\mathrm{rad}}$ rel. $L_1$ & 5.93190790e-03 & 2.50000000e-02 & {\color{green!60!black}PASS} \\
\bottomrule
\end{tabular}
\end{table}

\subsection*{Plots}
\textit{Plots not available --- run the test suite and the plot generator first.}

\newpage

\section{Mach~2 Radiative Shock -- Multigroup}
\label{sec:mach2_multigroup}

\subsection*{Description}
The same Mach~2 radiative shock as the gray diffusion test \cite{lowrie2008,lowrie1999}, but using multigroup radiation transport with 32 logarithmically spaced energy groups. This tests the frequency-dependent radiation solver and verifies that the multigroup solution reproduces the correct spatial profiles. In addition, the radiation spectrum of the hottest cell is compared to a Planck distribution at the local gas temperature to check spectral equilibrium.

\textbf{Code and physics aspects verified:}
\begin{itemize}
  \item \textbf{Multigroup radiation transport:} Unlike the gray test, the radiation field is resolved into 32 frequency groups. Each group has its own diffusion equation with frequency-dependent opacities. The test verifies that the group-by-group solver, the opacity evaluation, and the energy-group coupling are all correct.
  \item \textbf{Spectral equilibrium:} In the post-shock region where radiation and matter are in thermal equilibrium, the radiation spectrum should follow a Planck distribution at the local gas temperature. Comparing the computed spectrum to the Planck function checks that the multigroup solver correctly thermalises the radiation field.
  \item \textbf{Consistency with gray solution:} The multigroup spatial profiles (density, $T_{\mathrm{gas}}$, $T_{\mathrm{rad}}$) should closely match the gray solution. Significant deviations would indicate a bug in the frequency-dependent opacity averaging or group-boundary treatment.
  \item \textbf{Planck and Rosseland mean opacities:} The effective gray opacity emerges from averaging the frequency-dependent values. The test implicitly checks that the Planck and Rosseland means are computed consistently.
\end{itemize}

\textbf{Importance:} Production radiation-hydrodynamics simulations almost always use multigroup (or full-spectrum) transport. The gray approximation is a useful simplification for verification, but real applications require frequency dependence to capture line-driven effects, spectral hardening, and non-Planckian radiation fields. This test bridges the gap between the simple gray benchmark and production-level multigroup capability, ensuring that the frequency-dependent machinery does not introduce errors beyond those already present in the gray solution.

\subsection*{Initial Conditions}
Identical to the gray diffusion case (1024-cell Cartesian mesh, same left/right states and opacities). The energy groups span $E_{\min} = 10^{-3}\,k_B\!\cdot\!200\;\mathrm{K}$ to $E_{\max} = 10^{3}\,k_B\!\cdot\!200\;\mathrm{K}$, logarithmically spaced.

\subsection*{Boundary Conditions}
Inflow boundaries at left and right.

\subsection*{Mesh Movement}
Eulerian (fixed mesh).

\subsection*{Execution}
MPI, 8~CPUs, submitted via SLURM (partition \texttt{bigrun}, exclusive). Built with \texttt{--energy\_groups\_num=32}.

\subsection*{Pass Criteria}
Same spatial profile thresholds as the gray case:
\begin{itemize}
  \item Density: relative $L_1 \le 0.025$.
  \item Temperature: relative $L_1 \le 0.025$.
\end{itemize}

\subsection*{Achieved Results}
\begin{table}[htbp]
\centering
\begin{tabular}{l r r c}
\toprule
\textbf{Metric} & \textbf{Achieved} & \textbf{Threshold} & \textbf{Status} \\
\midrule
Density rel. $L_1$ & 3.13364936e-03 & 2.50000000e-02 & {\color{green!60!black}PASS} \\
Temperature rel. $L_1$ & 6.18345190e-03 & 2.50000000e-02 & {\color{green!60!black}PASS} \\
$T_{\mathrm{rad}}$ rel. $L_1$ & 5.94798497e-03 & 2.50000000e-02 & {\color{green!60!black}PASS} \\
\bottomrule
\end{tabular}
\end{table}

\subsection*{Plots}
\textit{Plots not available --- run the test suite and the plot generator first.}

\newpage

\section{Marshak Wave Problem~1 (Non-Equilibrium, Uniform Density)}
\label{sec:marshak_wave_1}

\subsection*{Description}
Problem~1 from Giron et~al.\ (2026, arXiv:2601.05120), originally defined as Test~2 in Krief \& McClarren (2024, arXiv:2401.05138). A supersonic non-equilibrium Marshak wave propagates into a cold, uniform-density medium. The opacity is strongly temperature-dependent: $\kappa_R = 100\,(T/\mathrm{keV})^{-3}$ and the Planck opacity is $\kappa_P = 0.001\,\kappa_R$ (far from equilibrium). The material energy is $u = 6.86\times10^{14}\,(T/\mathrm{keV})^4$.

\textbf{Code and physics aspects verified:}
\begin{itemize}
  \item \textbf{Grey flux-limited diffusion (FLD):} Exercises the implicit diffusion solver with power-law opacities and no flux limiter.
  \item \textbf{Non-equilibrium radiation--matter coupling:} With $\kappa_P \ll \kappa_R$, the radiation and gas temperatures decouple significantly, testing the absorption/emission coupling.
  \item \textbf{Time-dependent Marshak boundary condition.}
\end{itemize}

\subsection*{Initial Conditions}
Domain $x \in [0,\,0.2]$ with 512 uniform cells. $\rho = 1$ (uniform), initially cold ($T \approx 0$). The bath temperature at $x=0$ is $T_{\mathrm{bath}}(t) = 1.008038\,(t/\mathrm{ns})^{1/3}\;\mathrm{keV}$.

\subsection*{Boundary Conditions}
Marshak (radiation inflow) at left, zero-flux at right.

\subsection*{Mesh Movement}
Eulerian (fixed mesh).

\subsection*{Execution}
Serial, 1~CPU, direct execution.

\subsection*{Pass Criteria}
Relative $L_1$ error for both $T_{\mathrm{gas}}$ and $T_{\mathrm{rad}}$ must be $\le 10^{-2}$, compared to the self-similar analytical solution.

\subsection*{Achieved Results}
\begin{table}[htbp]
\centering
\begin{tabular}{l r r c}
\toprule
\textbf{Metric} & \textbf{Achieved} & \textbf{Threshold} & \textbf{Status} \\
\midrule
$T_{\mathrm{gas}}$ rel. $L_1$ & 6.246750e-03 & 1e-2 & {\color{green!60!black}PASS} \\
$T_{\mathrm{rad}}$ rel. $L_1$ & 6.279556e-03 & 1e-2 & {\color{green!60!black}PASS} \\
\bottomrule
\end{tabular}
\end{table}

\subsection*{Plots}
\textit{Plots not available --- run the test suite and the plot generator first.}

\newpage

\section{Marshak Wave Problem~2 (Equilibrium Limit, Uniform Density)}
\label{sec:marshak_wave_2}

\subsection*{Description}
Problem~2 from Giron et~al.\ (2026), originally Test~3 in Krief \& McClarren (2024). Same as Problem~1 except the Planck opacity equals the Rosseland opacity ($\kappa_P = \kappa_R$), placing the system in the equilibrium (LTE) limit where $T_{\mathrm{rad}} \approx T_{\mathrm{gas}}$.

\textbf{Code and physics aspects verified:}
\begin{itemize}
  \item \textbf{Equilibrium radiation diffusion:} Tests that the code correctly recovers the LTE limit where matter and radiation temperatures converge.
  \item \textbf{Marshak boundary with time-dependent driving temperature.}
\end{itemize}

\subsection*{Initial Conditions}
Domain $x \in [0,\,0.2]$ with 512 uniform cells, $\rho=1$, initially cold. $T_{\mathrm{bath}}(t) = 1.014565\,(t/\mathrm{ns})^{1/3}\;\mathrm{keV}$.

\subsection*{Boundary Conditions}
Marshak (radiation inflow) at left, zero-flux at right.

\subsection*{Mesh Movement}
Eulerian (fixed mesh).

\subsection*{Execution}
Serial, 1~CPU, direct execution.

\subsection*{Pass Criteria}
Relative $L_1$ error for both $T_{\mathrm{gas}}$ and $T_{\mathrm{rad}}$ must be $\le 10^{-2}$.

\subsection*{Achieved Results}
\begin{table}[htbp]
\centering
\begin{tabular}{l r r c}
\toprule
\textbf{Metric} & \textbf{Achieved} & \textbf{Threshold} & \textbf{Status} \\
\midrule
$T_{\mathrm{gas}}$ rel. $L_1$ & 5.017600e-03 & 1e-2 & {\color{green!60!black}PASS} \\
$T_{\mathrm{rad}}$ rel. $L_1$ & 9.875044e-03 & 1e-2 & {\color{green!60!black}PASS} \\
\bottomrule
\end{tabular}
\end{table}

\subsection*{Plots}
\textit{Plots not available --- run the test suite and the plot generator first.}

\newpage

\section{Marshak Wave Problem~3 (Non-Uniform Density, Power-Law Profile)}
\label{sec:marshak_wave_3}

\subsection*{Description}
Problem~3 from Giron et~al.\ (2026), originally Test~1 in Derei et~al.\ (2024, arXiv:2411.14891). The material density varies as $\rho(x) = x^{20/19}$, and both the opacity and the energy depend on density: $\kappa_R = 40\,(T/\mathrm{keV})^{-1.5}\,\rho^{1.2}$, $\kappa_P = 0.0025\,\kappa_R$, $u = 10^{14}\,(T/\mathrm{keV})^{3.4}\,\rho^{0.86}$.

\textbf{Code and physics aspects verified:}
\begin{itemize}
  \item \textbf{Density-dependent opacities and EOS:} Tests the code's ability to handle coupled density--temperature dependence in the diffusion coefficients.
  \item \textbf{Non-uniform initial density profile.}
\end{itemize}

\subsection*{Initial Conditions}
Domain $x \in [0,\,1]$, 512 uniform cells. $\rho(x) = x^{20/19}$, initially cold. $T_{\mathrm{bath}}(t) = 1.0470478\,(t/\mathrm{ns})^{86/57}\;\mathrm{keV}$.

\subsection*{Boundary Conditions}
Marshak at left, zero-flux at right.

\subsection*{Mesh Movement}
Eulerian (fixed mesh).

\subsection*{Execution}
Serial, 1~CPU, direct execution.

\subsection*{Pass Criteria}
Relative $L_1$ error for both $T_{\mathrm{gas}}$ and $T_{\mathrm{rad}}$ must be $\le 10^{-2}$.

\subsection*{Achieved Results}
\begin{table}[htbp]
\centering
\begin{tabular}{l r r c}
\toprule
\textbf{Metric} & \textbf{Achieved} & \textbf{Threshold} & \textbf{Status} \\
\midrule
$T_{\mathrm{gas}}$ rel. $L_1$ & 8.431773e-03 & 1e-2 & {\color{green!60!black}PASS} \\
$T_{\mathrm{rad}}$ rel. $L_1$ & 7.447134e-03 & 1e-2 & {\color{green!60!black}PASS} \\
\bottomrule
\end{tabular}
\end{table}

\subsection*{Plots}
\textit{Plots not available --- run the test suite and the plot generator first.}

\newpage

\section{Marshak Wave Problem~4 (Divergent Density, Stretched Grid)}
\label{sec:marshak_wave_4}

\subsection*{Description}
Problem~4 from Giron et~al.\ (2026), originally Test~3 in Derei et~al.\ (2024). The density diverges toward the origin as $\rho(x) = x^{-40/139}$. A stretched grid (geometric progression) is used to resolve the steep gradients near $x=0$. $\kappa_R = 2\,(T/\mathrm{keV})^{-4.5}\,\rho^{1.9}$, $\kappa_P = 5{\times}10^{-4}\,\kappa_R$, $u = 10^{14}\,(T/\mathrm{keV})^{6}\,\rho^{0.7}$.

\textbf{Code and physics aspects verified:}
\begin{itemize}
  \item \textbf{Stretched (non-uniform) grid:} The geometric progression mesh tests the diffusion solver on highly non-uniform cell sizes.
  \item \textbf{Divergent density at origin:} Stresses the opacity and energy evaluations near singular density profiles.
\end{itemize}

\subsection*{Initial Conditions}
Domain with 512 stretched cells (geometric progression). $\rho(x) = x^{-40/139}$, initially cold. $T_{\mathrm{bath}}(t) = 1.01008116\,(t/\mathrm{ns})^{14/139}\;\mathrm{keV}$.

\subsection*{Boundary Conditions}
Marshak at left, zero-flux at right.

\subsection*{Mesh Movement}
Eulerian (fixed mesh).

\subsection*{Execution}
Serial, 1~CPU, direct execution.

\subsection*{Pass Criteria}
Relative $L_1$ error for both $T_{\mathrm{gas}}$ and $T_{\mathrm{rad}}$ must be $\le 10^{-2}$.

\subsection*{Achieved Results}
\begin{table}[htbp]
\centering
\begin{tabular}{l r r c}
\toprule
\textbf{Metric} & \textbf{Achieved} & \textbf{Threshold} & \textbf{Status} \\
\midrule
$T_{\mathrm{gas}}$ rel. $L_1$ & 5.555073e-03 & 1e-2 & {\color{green!60!black}PASS} \\
$T_{\mathrm{rad}}$ rel. $L_1$ & 4.322846e-03 & 1e-2 & {\color{green!60!black}PASS} \\
\bottomrule
\end{tabular}
\end{table}

\subsection*{Plots}
\textit{Plots not available --- run the test suite and the plot generator first.}

\newpage

\section{Gresho Vortex (Eulerian)}
\label{sec:gresho_euler}

\subsection*{Description}
The Gresho vortex \cite{liska2003} is a stationary azimuthal flow in centrifugal equilibrium: the pressure gradient exactly balances the centripetal acceleration, so the exact solution at any time is the initial condition. Numerical dissipation and dispersion cause the vortex to decay over time, making this an excellent test of a code's ability to preserve rotating structures.

This test uses an \textbf{Eulerian (fixed) mesh}, which subjects the solution to advection errors as the vortex sweeps across the grid. The 3D setup uses a single cell in $z$ to approximate a 2D problem.

\textbf{Code and physics aspects verified:}
\begin{itemize}
  \item \textbf{Euler fluxes on a fixed mesh:} The vortex tests whether the Riemann solver and spatial reconstruction preserve a smooth, subsonic rotating flow.
  \item \textbf{Numerical dissipation:} The rate of vortex decay is a direct measure of the scheme's numerical viscosity.
\end{itemize}

\subsection*{Initial Conditions}
Domain $[-0.5,\,0.5]^2 \times [0,\,\Delta z]$, Cartesian $50\times50\times1$. $\gamma = 5/3$, $\rho = 1$. Azimuthal velocity: $v_\theta = 5r$ for $r<0.2$, $v_\theta = 2-5r$ for $0.2 \le r \le 0.4$, $v_\theta = 0$ for $r>0.4$. Pressure set for centrifugal equilibrium. Simulation runs to $t = 5$.

\subsection*{Boundary Conditions}
Rigid (reflective) walls on all faces.

\subsection*{Mesh Movement}
Eulerian (fixed mesh).

\subsection*{Execution}
Serial, 1~CPU, direct execution.

\subsection*{Pass Criteria}
Relative $L_1$ error for the azimuthal velocity profile (volume-averaged, radially binned) at $t=5$ compared to the initial condition must be $\le 0.1$.

\subsection*{Achieved Results}
\begin{table}[htbp]
\centering
\begin{tabular}{l r r c}
\toprule
\textbf{Metric} & \textbf{Achieved} & \textbf{Threshold} & \textbf{Status} \\
\midrule
$v_\theta$ rel. $L_1$ & 3.226244e-02 & 0.1 & {\color{green!60!black}PASS} \\
\bottomrule
\end{tabular}
\end{table}

\subsection*{Plots}
\textit{Plots not available --- run the test suite and the plot generator first.}

\newpage

\section{Gresho Vortex (Lagrangian + RoundCells)}
\label{sec:gresho_lagrangian}

\subsection*{Description}
Same Gresho vortex problem as the Eulerian case, but run with a \textbf{Lagrangian mesh} augmented by the Arepo-style cell-rounding correction (\texttt{RoundCells3D}). Mesh point motion is restricted to the $xy$-plane.

Because the mesh co-rotates with the fluid, advection errors are eliminated. The primary error source becomes the mesh deformation and the cell-rounding correction. The Lagrangian scheme should preserve the vortex significantly better than the Eulerian scheme.

\textbf{Code and physics aspects verified:}
\begin{itemize}
  \item \textbf{Lagrangian mesh motion:} Verifies that the moving mesh correctly tracks the rotating flow.
  \item \textbf{Cell rounding:} The RoundCells correction must maintain mesh quality without introducing excessive dissipation.
  \item \textbf{Custom point-motion wrapper:} The $z$-velocity zeroing wrapper is tested for correctness.
\end{itemize}

\subsection*{Initial Conditions}
Identical to the Eulerian case: $50\times50\times1$ Cartesian mesh, $\gamma = 5/3$, same velocity and pressure profiles, $t_{\mathrm{end}} = 5$.

\subsection*{Boundary Conditions}
Rigid (reflective) walls on all faces.

\subsection*{Mesh Movement}
Lagrangian + RoundCells (restricted to $xy$-plane).

\subsection*{Execution}
MPI, 8~CPUs, submitted via SLURM (partition \texttt{bigrun}, exclusive).

\subsection*{Pass Criteria}
Relative $L_1$ error for the azimuthal velocity profile must be $\le 0.05$.

\subsection*{Achieved Results}
\begin{table}[htbp]
\centering
\begin{tabular}{l r r c}
\toprule
\textbf{Metric} & \textbf{Achieved} & \textbf{Threshold} & \textbf{Status} \\
\midrule
$v_\theta$ rel. $L_1$ & 3.196482e-02 & 0.05 & {\color{green!60!black}PASS} \\
\bottomrule
\end{tabular}
\end{table}

\subsection*{Plots}
\textit{Plots not available --- run the test suite and the plot generator first.}

\newpage

\section{Spherical LSQ Gradient -- Linear Field Verification}
\label{sec:spherical_gauss_linear}

\subsection*{Description}
A pure-reconstruction test that verifies the weighted least-squares (LSQ) gradient in spherical coordinate space implemented in \texttt{SphericalLinearGauss3D}. Fields that are exactly linear in $(r,\,\theta)$ are prescribed at every cell centre and then reconstructed to face centroids. Because the LSQ gradient can represent linear functions exactly, the interpolated values must agree with the exact values to machine precision.

\textbf{Code and physics aspects verified:}
\begin{itemize}
  \item \textbf{LSQ gradient accuracy:} The 3\texttimes3 normal system and its analytic inverse must yield the exact coordinate derivatives $\partial f/\partial r$, $\partial f/\partial\theta$, $\partial f/\partial\phi$.
  \item \textbf{Velocity frame conversion:} Each cell's velocity is projected onto its own local spherical basis; the interpolated velocity at the face centroid is converted back to Cartesian using the basis at that point.
  \item \textbf{Coordinate-space interpolation:} The displacement $(\Delta r,\,\Delta\theta,\,\Delta\phi)$ between cell centre and face centroid must be computed correctly.
\end{itemize}

\subsection*{Initial Conditions}
A Voronoi tessellation from $\sim$4\,000 random points in a spherical shell $1 < r < 2$ embedded in a $[-4,\,4]^3$ box. Scalar fields (density, internal energy) are linear in $(r,\,\theta)$; velocity components $(v_r,\,v_\theta,\,v_\phi)$ are also linear in $(r,\,\theta)$ (no $\phi$ dependence to avoid angle-wrapping artefacts). The slope limiter is disabled so that the reconstruction is purely linear.

\subsection*{Boundary Conditions}
Rigid walls at box faces.

\subsection*{Mesh Movement}
Static (no time evolution).

\subsection*{Execution}
Serial, 1~CPU, direct execution.

\subsection*{Pass Criteria}
For interior faces (both cell CMs in $1.1 < r < 1.9$):
\begin{itemize}
  \item Scalar max relative error $< 10^{-8}$.
  \item Velocity max relative error $< 0.1$.
\end{itemize}

\subsection*{Achieved Results}
\begin{table}[htbp]
\centering
\begin{tabular}{l r r c}
\toprule
\textbf{Metric} & \textbf{Achieved} & \textbf{Threshold} & \textbf{Status} \\
\midrule
Scalar max rel. error & 5.658526829483e-16 & 1e-08 & {\color{green!60!black}PASS} \\
Velocity max rel. error & 5.568379703788e-12 & 1e-01 & {\color{green!60!black}PASS} \\
Faces checked & 87600 & > 0 & {\color{green!60!black}PASS} \\
\bottomrule
\end{tabular}
\end{table}

\subsection*{Plots}
\textit{This test does not produce comparison plots.}

\newpage

\section{Spherical Shell Collapse (Symmetry Test)}
\label{sec:spherical_collapse}

\subsection*{Description}
A dense spherical shell ($0.9 < r < 1.0$, $\rho=1$, $P=1$) is embedded in a low-density ambient medium ($\rho=10^{-3}$, $P=10^{-4}$) and allowed to collapse inward on a fixed Eulerian mesh. The mesh is constructed by generating a well-rounded template sphere (\texttt{RandSphereSurfaceRounded}, $\sim$10\,000 points) and replicating it at logarithmically spaced radii from $R=1.1$ down to $R=0.05$ with constant $\Delta R/R$, ensuring near-unity aspect ratio. The inner core ($r<0.05$) and the region between the outer sphere and the box walls are filled with random points.

The simulation runs until the volume-averaged inward velocity at $r=0.05$ reaches unity. Snapshots are written each time the shock front advances inward by $0.1$ in radius.

\textbf{Code and physics aspects verified:}
\begin{itemize}
  \item \textbf{Spherical symmetry preservation:} The primary metric is the angular scatter (standard deviation / mean) of density and radial velocity in each radial bin. A well-behaved code should maintain low scatter throughout the collapse.
  \item \textbf{Eulerian advection on unstructured mesh:} The fixed mesh forces all dynamics to be captured by inter-cell fluxes, stressing the Riemann solver and reconstruction.
  \item \textbf{Strong shock convergence:} The collapsing shell drives a converging shock that amplifies any asymmetry.
\end{itemize}

\subsection*{Initial Conditions}
Dense shell: $\rho = 1$, $P = 1$ for $0.9 < r < 1.0$; ambient: $\rho = 10^{-3}$, $P = 10^{-4}$ elsewhere. Zero velocity everywhere. $\gamma = 5/3$.

\subsection*{Boundary Conditions}
Rigid (reflective) walls on all faces of the $[-1.55, 1.55]^3$ box.

\subsection*{Mesh Movement}
Eulerian (fixed mesh).

\subsection*{Execution}
MPI, 64~CPUs, submitted via SLURM (partition \texttt{bigrun}, exclusive).

\subsection*{Pass Criteria}
Maximum angular density scatter $< 0.1$ and maximum angular velocity scatter $< 0.1$ across all radial bins with significant density ($\rho > 0.01$) or velocity ($|v_r| > 0.01$).

\subsection*{Achieved Results}
\begin{table}[htbp]
\centering
\begin{tabular}{l r r c}
\toprule
\textbf{Metric} & \textbf{Achieved} & \textbf{Threshold} & \textbf{Status} \\
\midrule
Max density scatter & 1.970876181564e-05 & 0.1 & {\color{green!60!black}PASS} \\
Max velocity scatter & 1.279011537029e-05 & 0.1 & {\color{green!60!black}PASS} \\
\bottomrule
\end{tabular}
\end{table}

\subsection*{Plots}
\textit{Plots not available --- run the test suite and the plot generator first.}

\newpage

\section{Rayleigh--Taylor Instability (3D, MPI)}
\label{sec:rayleigh_taylor_mpi}

\subsection*{Description}
The Rayleigh--Taylor instability \cite{rayleigh1883,taylor1950rt} develops when a heavy fluid is supported against gravity above a lighter fluid. A small sinusoidal perturbation at the interface grows exponentially in the linear regime with a growth rate determined by the Atwood number, gravitational acceleration, and perturbation wavenumber \cite{chandrasekhar1961}.

The initial condition places $\rho_{\mathrm{heavy}}=2$ above $\rho_{\mathrm{light}}=1$ in a $[0,1]^2 \times [0,2]$ box with a flat interface at $z=1$ in exact hydrostatic equilibrium. A two-mode cosine velocity perturbation in $v_z$ (amplitude $0.03$, Gaussian-localised near the interface) seeds the instability. Constant downward gravity $g=0.5$ is applied.

\textbf{Code and physics aspects verified:}
\begin{itemize}
  \item \textbf{Source term coupling:} The conservative gravity source term (\texttt{ConservativeForce3D}) must correctly inject momentum and energy to maintain hydrostatic equilibrium and drive the instability.
  \item \textbf{Lagrangian mesh motion with RoundCells:} The moving mesh must track the developing interface while keeping cells well-shaped.
  \item \textbf{MPI scalability:} The test runs on 128~MPI tasks with $\sim$10$^6$ cells, exercising the parallel domain decomposition and communication.
\end{itemize}

\subsection*{Initial Conditions}
$\rho_{\mathrm{light}}=1$ for $z < 1$, $\rho_{\mathrm{heavy}}=2$ for $z \geq 1$ (flat interface). Hydrostatic pressure with $P_0=10$, $g=0.5$. $\gamma=5/3$. Velocity perturbation: $v_z = 0.03\,[\cos(2\pi x) + \cos(2\pi y)]\,\exp[-(z-1)^2/0.04]$.

\subsection*{Boundary Conditions}
Rigid (reflective) walls on all faces of the $[0,1]^2 \times [0,2]$ box.

\subsection*{Mesh Movement}
Lagrangian + RoundCells (\texttt{Lagrangian3D} wrapped in \texttt{RoundCells3D}).

\subsection*{Execution}
MPI, 128~CPUs, submitted via SLURM (partition \texttt{bigrun}, exclusive).

\subsection*{Pass Criteria}
Fitted linear growth rate $\sigma_{\mathrm{fit}}$ (from $t=2$ to $t=3$) within 25\% of analytical $\sigma = \sqrt{A\,g\,k}$ where $A=1/3$, $g=0.5$, $k=2\pi$.

\subsection*{Achieved Results}
\begin{table}[htbp]
\centering
\begin{tabular}{l r r c}
\toprule
\textbf{Metric} & \textbf{Achieved} & \textbf{Threshold} & \textbf{Status} \\
\midrule
Growth rate rel. error & 9.56195392e-02 & 2.50000000e-01 & {\color{green!60!black}PASS} \\
Fitted $\sigma$ & 9.25476680e-01 & (analytical: 1.02332671e+00) & {\color{green!60!black}PASS} \\
\bottomrule
\end{tabular}
\end{table}

\subsection*{Plots}
\textit{Plots not available --- run the test suite and the plot generator first.}

\newpage

\section{1D Eulerian Diffusion Suite (Gray Free--Free, Cooling Limiter)}
\label{sec:eulerian_diffusion_freefree_suite}

\subsection*{Description}
A four-case resolution and cooling-limiter study for 1D Eulerian radiation-hydrodynamics with gray flux-limited diffusion, free--free (bremsstrahlung) absorption opacity, and Thomson scattering. Two counter-propagating streams ($v_x = \pm 10^8\,\mathrm{cm\,s^{-1}}$) collide at the domain midpoint, driving a radiative shock. The suite runs the same physics at two resolutions (512 and 32~cells) and with the cooling-time opacity limiter toggled on and off, producing a four-way comparison.

\textbf{The cooling-time opacity limiter.}\enspace When the post-shock cooling length is under-resolved, the finite-volume scheme cannot capture the thin radiative relaxation layer behind the shock. Because energy is radiated away over fewer cells than the physics requires, the post-shock temperature in an under-resolved run is artificially elevated: the gas cools too fast relative to the compression time, violating the resolution-independent structure of the shock.

The limiter addresses this by comparing two time scales in every cell that is undergoing compression:
\begin{enumerate}
  \item \textbf{Hydrodynamic heating time:} $t_{\mathrm{hydro}} = 1/\max(-\nabla\!\cdot\!\mathbf{v},\,\epsilon)$, measuring how fast compression is depositing energy.
  \item \textbf{Radiative cooling time:} $t_{\mathrm{cool}} = \rho\,e / (P_{\mathrm{Planck}} + P_{\mathrm{Compton}})$, where $P_{\mathrm{Planck}} = c\,\sigma_P\,(aT^4 - E_r)$ and $P_{\mathrm{Compton}}$ is the Compton exchange rate (if Compton scattering is enabled).
\end{enumerate}
The limiter activates only in cells where the compression speed exceeds a fraction of the bulk velocity ($(-\nabla\!\cdot\!\mathbf{v})\,\Delta x > 0.25\,|\mathbf{v}|$). If $t_{\mathrm{cool}} < 2\,t_{\mathrm{hydro}}$, both the Planck absorption and scattering opacities are scaled down by a common factor so that the effective cooling time equals $2\,t_{\mathrm{hydro}}$. This preserves the relative weighting of absorption and scattering while preventing the cooling rate from exceeding what the grid can resolve.

\textbf{Code and physics aspects verified:}
\begin{itemize}
  \item \textbf{Eulerian radiation-hydrodynamics coupling:} Tests the full hydro + gray diffusion pipeline in a shock-forming setup with counter-propagating inflows.
  \item \textbf{Free--free opacity + Thomson scattering:} Validates the gray Kramers free--free absorption and constant Thomson scattering opacity paths.
  \item \textbf{Resolution convergence:} Comparing 512 vs.\ 32~cells reveals the sensitivity of the post-shock structure to resolution. The high-resolution run serves as a reference.
  \item \textbf{Cooling-limiter effectiveness:} With the limiter enabled, the 32-cell $T_{\mathrm{gas}}$ peak should be closer to the 512-cell result, demonstrating that the limiter successfully prevents resolution-dependent over-cooling.
  \item \textbf{Mixed radiation boundary conditions:} The radiation boundary is open on the $x$-faces and closed on $y/z$, testing a directional boundary implementation.
\end{itemize}

\textbf{Importance:} Under-resolved cooling layers are one of the most common failure modes in production radiation-hydrodynamics simulations. This suite directly tests the optional limiter that mitigates this problem and quantifies its effect across an order-of-magnitude resolution difference.

\subsection*{Initial Conditions}
Domain $x \in [0,\,2\times10^{12}]\,\mathrm{cm}$ with uniform cells. The gas is initialised with $\rho = 2\times10^{-13}\,\mathrm{g\,cm^{-3}}$, $T = 2\times10^{5}\,\mathrm{K}$. The left half has $v_x = +10^8\,\mathrm{cm\,s^{-1}}$ and the right half $v_x = -10^8\,\mathrm{cm\,s^{-1}}$, forming a symmetric collision. The adiabatic index is $\gamma = 5/3$.

The four runs are:
\begin{itemize}
  \item 512 cells, limiter \textbf{off} (reference).
  \item 512 cells, limiter \textbf{on}.
  \item 32 cells, limiter \textbf{off}.
  \item 32 cells, limiter \textbf{on}.
\end{itemize}
The run terminates when the shock front (maximum $x$ with $\rho > 2\rho_0$) reaches $0.75\,L$, or when $t = 9\times10^4\,\mathrm{s}$.

\subsection*{Boundary Conditions}
Hydrodynamics: left/right inflow ghost states matching the respective initial half-domain states; rigid walls on $y/z$ faces. Radiation: open (streaming) on the two $x$-boundary faces, closed on $y/z$.

\subsection*{Mesh Movement}
Eulerian (fixed mesh).

\subsection*{Execution}
MPI, 16~CPUs for 512-cell runs, 4~CPUs for 32-cell runs, submitted via SLURM (partition \texttt{bigrun}, exclusive).

\subsection*{Pass Criteria}
Each of the four sub-runs must complete without fatal errors and produce a valid \texttt{temperature\_profile.txt} with at least two data points. The suite then generates four-way comparison plots (gas temperature, radiation temperature, density, and velocity vs.\ $x$) which are inspected visually. Quantitative acceptance:
\begin{itemize}
  \item All profile files contain finite, positive temperatures.
  \item The 32-cell limited peak $T_{\mathrm{gas}}$ is closer to the 512-cell reference than the 32-cell unlimited peak.
\end{itemize}

\subsection*{Achieved Results}
\begin{table}[htbp]
\centering
\begin{tabular}{l r r c}
\toprule
\textbf{Metric} & \textbf{Achieved} & \textbf{Threshold} & \textbf{Status} \\
\midrule
512: points & 512 & $\ge 2$ & {\color{green!60!black}PASS} \\
512: peak $T_{\mathrm{gas}}$ & 1.9532 keV & finite & {\color{green!60!black}PASS} \\
512 limited: points & 512 & $\ge 2$ & {\color{green!60!black}PASS} \\
512 limited: peak $T_{\mathrm{gas}}$ & 1.9314 keV & finite & {\color{green!60!black}PASS} \\
32: points & 32 & $\ge 2$ & {\color{green!60!black}PASS} \\
32: peak $T_{\mathrm{gas}}$ & 0.0142 keV & finite & {\color{green!60!black}PASS} \\
32 limited: points & 32 & $\ge 2$ & {\color{green!60!black}PASS} \\
32 limited: peak $T_{\mathrm{gas}}$ & 1.3024 keV & finite & {\color{green!60!black}PASS} \\
\bottomrule
\end{tabular}
\end{table}

\subsection*{Plots}
\begin{figure}[htbp]
  \centering
  \includegraphics[width=0.48\textwidth]{cases/eulerian_diffusion_freefree_compare/temperature_vs_x_compare_512_512_limited_32_32_limited.pdf}
  \hfill
  \includegraphics[width=0.48\textwidth]{cases/eulerian_diffusion_freefree_compare/trad_vs_x_compare_512_512_limited_32_32_limited.pdf}
  \\[6pt]
  \includegraphics[width=0.48\textwidth]{cases/eulerian_diffusion_freefree_compare/density_vs_x_compare_512_512_limited_32_32_limited.pdf}
  \hfill
  \includegraphics[width=0.48\textwidth]{cases/eulerian_diffusion_freefree_compare/velocity_vs_x_compare_512_512_limited_32_32_limited.pdf}
  \caption{Gray free--free diffusion suite: four-way comparison of 512-cell (solid), 512-limited (dash-dot), 32-cell (dashed), and 32-limited (dotted) runs. Top row: gas temperature (left) and radiation temperature (right). Bottom row: density (left) and $x$-velocity (right). The cooling limiter brings the low-resolution peak $T_{\mathrm{gas}}$ closer to the high-resolution reference.}
  \label{fig:eulerian_diffusion_freefree_suite}
\end{figure}

\newpage

\section{1D Eulerian Diffusion Suite (Multigroup Free--Free, Cooling Limiter)}
\label{sec:eulerian_diffusion_freefree_multigroup_suite}

\subsection*{Description}
The multigroup counterpart of the gray free--free suite. The same symmetric collision problem is run with 32 logarithmically spaced energy groups, frequency-dependent free--free absorption, and Compton scattering. The same four-variant pattern is used (512/32~cells $\times$ limiter on/off) and a four-way comparison of profiles is produced.

\textbf{The multigroup cooling limiter.}\enspace In multigroup mode the cooling-time estimate sums over all energy groups: the Planck exchange is accumulated group-by-group as $c\,\sigma_{P,g}\,(B_g(T) - E_g)$ and the Compton exchange is computed from the full Compton transfer matrix $S_{g,g'}$ rather than a gray approximation. When the net cooling time is shorter than $2\,t_{\mathrm{hydro}}$, all group-level absorption opacities are scaled by a common factor and the Compton coupling matrix entries ($S$ and $\mathrm{d}S/\mathrm{d}U_m$) are multiplied by a stored per-cell scale factor, ensuring that the full multigroup exchange is consistently limited.

\textbf{Code and physics aspects verified:}
\begin{itemize}
  \item \textbf{Multigroup diffusion with frequency-dependent opacities:} Each of the 32 energy groups evolves with its own absorption coefficient, testing the group-by-group solver.
  \item \textbf{Compton scattering matrix:} The full Kompaneets-based Compton transfer matrix is exercised in a dynamic (non-static) setting with strong temperature gradients.
  \item \textbf{Multigroup cooling limiter:} Verifies that the per-cell Compton-limiter scale factor is correctly stored, applied to the $S$ and $\mathrm{d}S/\mathrm{d}U_m$ matrices, and produces the same qualitative resolution-convergence improvement seen in the gray case.
  \item \textbf{Consistency with gray results:} The spatial profiles should be qualitatively similar to the gray suite, with frequency-dependent effects introducing only modest quantitative differences.
\end{itemize}

\textbf{Importance:} Production simulations typically use multigroup transport and Compton scattering. This suite ensures that the cooling limiter---originally developed for the gray solver---extends correctly to the frequency-dependent case including the Compton coupling terms.

\subsection*{Initial Conditions}
Identical domain and gas state to the gray suite ($L = 2\times10^{12}\,\mathrm{cm}$, $\rho = 2\times10^{-13}\,\mathrm{g\,cm^{-3}}$, $T = 2\times10^{5}\,\mathrm{K}$, symmetric $\pm 10^8\,\mathrm{cm\,s^{-1}}$). Radiation transport uses 32 logarithmically spaced energy groups with free--free opacity (frequency-dependent Kramers formula) and Compton scattering.

The four runs are:
\begin{itemize}
  \item 512 cells, limiter \textbf{off} (reference).
  \item 512 cells, limiter \textbf{on}.
  \item 32 cells, limiter \textbf{off}.
  \item 32 cells, limiter \textbf{on}.
\end{itemize}

\subsection*{Boundary Conditions}
Same as the gray suite: inflow on $x$, rigid on $y/z$; radiation open on $x$, closed on $y/z$.

\subsection*{Mesh Movement}
Eulerian (fixed mesh).

\subsection*{Execution}
MPI, 16~CPUs for 512-cell runs, 4~CPUs for 32-cell runs, submitted via SLURM (partition \texttt{bigrun}, exclusive). Built with \texttt{--energy\_groups\_num=32}.

\subsection*{Pass Criteria}
Same structural criteria as the gray suite: all four sub-runs must complete and produce valid profiles. Additionally:
\begin{itemize}
  \item All profile files contain finite, positive temperatures.
  \item The 32-cell limited peak $T_{\mathrm{gas}}$ is closer to the 512-cell reference than the 32-cell unlimited peak.
\end{itemize}

\subsection*{Achieved Results}
\begin{table}[htbp]
\centering
\begin{tabular}{l r r c}
\toprule
\textbf{Metric} & \textbf{Achieved} & \textbf{Threshold} & \textbf{Status} \\
\midrule
512: points & 512 & $\ge 2$ & {\color{green!60!black}PASS} \\
512: peak $T_{\mathrm{gas}}$ & 0.6378 keV & finite & {\color{green!60!black}PASS} \\
512 limited: points & 512 & $\ge 2$ & {\color{green!60!black}PASS} \\
512 limited: peak $T_{\mathrm{gas}}$ & 0.8820 keV & finite & {\color{green!60!black}PASS} \\
32: points & 32 & $\ge 2$ & {\color{green!60!black}PASS} \\
32: peak $T_{\mathrm{gas}}$ & 0.0128 keV & finite & {\color{green!60!black}PASS} \\
32 limited: points & 32 & $\ge 2$ & {\color{green!60!black}PASS} \\
32 limited: peak $T_{\mathrm{gas}}$ & 0.1429 keV & finite & {\color{green!60!black}PASS} \\
\bottomrule
\end{tabular}
\end{table}

\subsection*{Plots}
\begin{figure}[htbp]
  \centering
  \includegraphics[width=0.48\textwidth]{cases/eulerian_diffusion_freefree_multigroup_compare/temperature_vs_x_compare_mg32_512_512_limited_32_32_limited.pdf}
  \hfill
  \includegraphics[width=0.48\textwidth]{cases/eulerian_diffusion_freefree_multigroup_compare/trad_vs_x_compare_mg32_512_512_limited_32_32_limited.pdf}
  \\[6pt]
  \includegraphics[width=0.48\textwidth]{cases/eulerian_diffusion_freefree_multigroup_compare/density_vs_x_compare_mg32_512_512_limited_32_32_limited.pdf}
  \hfill
  \includegraphics[width=0.48\textwidth]{cases/eulerian_diffusion_freefree_multigroup_compare/velocity_vs_x_compare_mg32_512_512_limited_32_32_limited.pdf}
  \caption{Multigroup free--free diffusion suite (32~energy groups): four-way comparison analogous to the gray suite. Top row: gas temperature (left) and radiation temperature (right). Bottom row: density (left) and $x$-velocity (right). The cooling limiter reduces the low-resolution temperature excess as in the gray case, while the multigroup Compton coupling introduces modest quantitative differences.}
  \label{fig:eulerian_diffusion_freefree_multigroup_suite}
\end{figure}

\newpage


\begin{thebibliography}{99}

% --- V&V methodology ---

\bibitem{roache1998}
P.~J. Roache,
\textit{Verification and Validation in Computational Science and Engineering},
Hermosa Publishers, Albuquerque, NM, 1998.

\bibitem{oberkampf2010}
W.~L. Oberkampf and C.~J. Roy,
\textit{Verification and Validation in Scientific Computing},
Cambridge University Press, 2010.

\bibitem{aiaa1998}
AIAA,
``Guide for the Verification and Validation of Computational Fluid Dynamics
Simulations,'' AIAA-G-077-1998, American Institute of Aeronautics and
Astronautics, 1998.

\bibitem{asme2006}
ASME,
``Guide for Verification and Validation in Computational Solid Mechanics,''
ASME V\&V~10-2006, American Society of Mechanical Engineers, 2006.

\bibitem{salari2000}
K.~Salari and P.~Knupp,
``Code Verification by the Method of Manufactured Solutions,''
Sandia Report SAND2000-1444, Sandia National Laboratories, 2000.

\bibitem{roache2002}
P.~J. Roache,
``Code Verification by the Method of Manufactured Solutions,''
\textit{J.~Fluids Eng.}, 124(1):4--10, 2002.

\bibitem{richardson1911}
L.~F. Richardson,
``The Approximate Arithmetical Solution by Finite Differences of Physical
Problems Involving Differential Equations, with an Application to the Stresses
in a Masonry Dam,''
\textit{Philos. Trans. R. Soc. London, Ser.~A}, 210:307--357, 1911.

% --- Sod shock tube ---

\bibitem{sod1978}
G.~A. Sod,
``A Survey of Several Finite Difference Methods for Systems of Nonlinear
Hyperbolic Conservation Laws,''
\textit{J.~Comput. Phys.}, 27(1):1--31, 1978.

\bibitem{toro2009}
E.~F. Toro,
\textit{Riemann Solvers and Numerical Methods for Fluid Dynamics: A Practical
Introduction}, 3rd~ed., Springer, 2009.

% --- Sedov--Taylor blast wave ---

\bibitem{sedov1959}
L.~I. Sedov,
\textit{Similarity and Dimensional Methods in Mechanics},
Academic Press, New York, 1959.

\bibitem{taylor1950}
G.~I. Taylor,
``The Formation of a Blast Wave by a Very Intense Explosion. I.~Theoretical
Discussion,''
\textit{Proc. R. Soc. London, Ser.~A}, 201(1065):159--174, 1950.

\bibitem{kamm2007}
J.~R. Kamm and F.~X. Timmes,
``On Efficient Generation of Numerically Robust Sedov Solutions,''
Los Alamos Report LA-UR-07-2849, 2007.

% --- Till Compton equilibration ---

\bibitem{till2021}
A.~Till,
``Exploring Compton Scattering Options for Thermal Radiative Transfer,''
in \textit{Proc. ANS M\&C 2021}, LA-UR-20-29430, 2021.

\bibitem{mcgraw2023}
C.~McGraw, A.~Till, and J.~Warsa,
``A New Operator-Split Compton Scattering Method,''
\textit{J.~Comput. Phys.}, 478:111980, 2023.

% --- Lane--Emden self-gravity ---

\bibitem{chandrasekhar1939}
S.~Chandrasekhar,
\textit{An Introduction to the Study of Stellar Structure},
University of Chicago Press, 1939.

\bibitem{lane1870}
J.~Homer Lane,
``On the Theoretical Temperature of the Sun, under the Hypothesis of a Gaseous
Mass Maintaining Its Volume by Its Internal Heat, and Depending on the Laws of
Gases as Known to Terrestrial Experiment,''
\textit{Am. J. Sci.}, 2nd Series, 50(148):57--74, 1870.

% --- Mach 2 radiative shock ---

\bibitem{lowrie2008}
R.~B. Lowrie and J.~D. Edwards,
``Radiative Shock Solutions with Grey Nonequilibrium Diffusion,''
\textit{Shock Waves}, 18(2):129--143, 2008.

\bibitem{lowrie1999}
R.~B. Lowrie, J.~E. Morel, and J.~A. Hittinger,
``The Coupling of Radiation and Hydrodynamics,''
\textit{Astrophys. J.}, 521(1):432--450, 1999.

\bibitem{mihalas1984}
D.~Mihalas and B.~W. Mihalas,
\textit{Foundations of Radiation Hydrodynamics},
Oxford University Press, 1984.

\bibitem{zel2002}
Ya.~B. Zel'dovich and Yu.~P. Raizer,
\textit{Physics of Shock Waves and High-Temperature Hydrodynamic Phenomena},
Dover Publications, 2002 (reprint).

% --- Gresho vortex ---

\bibitem{liska2003}
R.~Liska and B.~Wendroff,
``Comparison of Several Difference Schemes on 1D and 2D Test Problems for the
Euler Equations,''
\textit{SIAM J.~Sci.~Comput.}, 25(3):995--1017, 2003.

% --- Rayleigh--Taylor instability ---

\bibitem{rayleigh1883}
Lord Rayleigh,
``Investigation of the Character of the Equilibrium of an Incompressible Heavy
Fluid of Variable Density,''
\textit{Proc.~London Math.~Soc.}, s1-14(1):170--177, 1883.

\bibitem{taylor1950rt}
G.~I.~Taylor,
``The Instability of Liquid Surfaces when Accelerated in a Direction
Perpendicular to their Planes.~I,''
\textit{Proc.~R.~Soc.~Lond.~A}, 201(1065):192--196, 1950.

\bibitem{chandrasekhar1961}
S.~Chandrasekhar,
\textit{Hydrodynamic and Hydromagnetic Stability},
Oxford University Press, 1961.

\end{thebibliography}


\end{document}
